\documentclass[12pt]{article}
\usepackage[utf8]{inputenc}
\usepackage[T1]{fontenc}
\usepackage{hyperref}
\title{Tonga Language Resources for K-12 Curriculum Development}
\author{Compiled by Claude Sonnet 5 (claude-sonnet-5) for Andr\'as Kornai}
\date{}
\begin{document}
\maketitle

\section{Introduction}
This dossier covers Tonga (ISO 639-3 toh; Glottocode gito1238), usually called Gitonga, of Inhambane Province, southern Mozambique. Speaker estimates vary widely, from about 327,000 to 553,000 depending on the source. It must not be confused with tog or toi. The search found Bible-related and audio materials but no dictionary, textbook or school materials.

\section{Primary Educational Resources}
None found in search results.

\section{Linguistic and Reference Materials}
No dictionary or grammar was found in the search. Wikipedia describes the language and reports closeness to Chopi.
\begin{itemize}
\item Gitonga language (Wikipedia mirror): \url{https://demo.azizisearch.com/lite/wikipedia/page/Gitonga_language}
\item Tonga language (Mozambique), Wikipedia: \url{https://wikipedia.com/wiki/ISO_639:toh}
\end{itemize}

\section{Religious and Community Resources}
\begin{itemize}
\item Article on Friar Amaral Bernardo Amaral's Bible translation into Gitonga (1984--2014): \url{https://unisapressjournals.co.za/index.php/OHJSA/article/view/13}
\item Report on the Christian Greek Scriptures released in Gitonga (2023): \url{https://pr418.com/en/news/region/mozambique/Christian-Greek-Scriptures-Released-in-Gitonga}
\item Joshua Project people group profile: \url{https://joshuaproject.net/people_groups/19123/MZ}
\end{itemize}

\section{Digital and Audio Resources}
Global Recordings Network lists Gitonga audio, reported to include the Jesus Film and a Faith Comes By Hearing New Testament recording: \url{https://glorecnet-website.k.globalrecordings.net/en/language/toh}

\section{Curriculum Development Context}
There is no evidence of school use. Written Gitonga exists mainly through Bible translations, which could serve as a text base. A dictionary, orthography and reading materials still have to be developed. Portuguese is the school language in Mozambique.
\end{document}